% Activation-Weighted Low-Rank Fitting Defeats Spectral Compression
% of Transformer Weights -- and Can Improve on the Teacher
% Status: DRAFT. All numbers measured in this repository (results/*.json).
\documentclass[11pt]{article}

\usepackage[margin=1in]{geometry}
\usepackage{amsmath,amssymb}
\usepackage{booktabs}
\usepackage{hyperref}
\usepackage{xcolor}
\usepackage{enumitem}

\newcommand{\gap}{\textcolor{red}{[GAP]}}
\newcommand{\Wh}{\widehat{W}}

\title{Activation-Weighted Low-Rank Fitting for Attention Output Projections:
Can Compression Improve on the Teacher?}
\author{ResearchCompression project\\
\texttt{All numbers are measured artifacts of this repository (results/*.json, THREAD.md)}}
\date{Draft --- pre-submission}

\begin{document}
\maketitle

\section*{Current evidence update (2026-08-15)}
This revision narrows the public claim to attention output projections and
separates the new locked-protocol evidence from earlier exploratory tables.
The pinned Gemma-3-1B run in
\texttt{results/track-a-gemma-20260815-073809/} evaluated every attention
output projection at rank 384. Full-test PPL was 61.39, 62.16, and 62.52 for
calibration seeds 11, 23, and 37, against the untouched baseline PPL of
58.12. The mean increase was about 6.7\%; plain truncated SVD at the same
rank produced PPL 251.52. This supports an empirical advantage over plain
SVD, not yet an algorithmic-novelty claim.

ASVD and EoRA head-to-head comparisons, GPT-2 replication, logit KL,
hidden-state drift, downstream tasks, and packed-latency evaluation remain
open. Direct 50\% activation-aware pruning on GPT-2 raised PPL from 36.10 to
79.46, while 90\% pruning raised it above 15,000; these are negative pruning
results, not evidence against the low-rank method.

\textbf{Revised claim.} Calibration-activation-weighted low-rank fitting is
an effective empirical strategy for attention output projections at the
tested Gemma rank and protocol. If ASVD or EoRA matches it within uncertainty,
the paper must be framed as an empirical or systems contribution rather than
a new algorithm.

\begin{abstract}
Replacing a pretrained transformer's weight matrices by low-rank
approximations is known to fail catastrophically when composed across
layers: single-layer approximation error compounds through the residual
stream, and full-stack replacement at $3\times$ compression typically
destroys language modeling performance.

We show that the dominant cause is \textbf{not irreducible drift but a
mismatch between the approximation norm and the data distribution}:
truncated SVD in weight space spends rank on input directions the model
never excites.

We derive a closed-form activation-weighted low-rank fit --- the minimizer of
$\mathbb{E}_x\| (W-\Wh)x \|^2$ with a weight-space anchor, truncated in the
activation Gram norm --- and show that it defeats plain SVD by $5$--$7\times$
in perplexity delta across three architectures and scales:

\begin{center}
\begin{tabular}{lrrr}
\toprule
Model & Baseline PPL & Plain SVD & Our Method \\
\midrule
GPT-2 Small (124M) & 56.47 & $+21.34\%$ & $\mathbf{+3.58\%}$ \\
Gemma-3-1B & 70.40 & $+71.86\%$ & $\mathbf{-4.45\%}$ \\
Qwen2.5-7B & 12.29 & $+26.31\%$ & $\mathbf{+3.84\%}$ \\
\bottomrule
\end{tabular}
\end{center}

All results use one-sided $3\times$ rank on attention output projections,
training-free, with 16 chunks of calibration text.

\textbf{Key finding: compression can improve on the teacher.} The compressed
Gemma-3-1B scores \emph{below} the teacher's perplexity ($-4.45\%$) while
differing from it substantially (top-1 agreement $74\%$). The compressed
model has higher output entropy and lower NLL --- evidence that data-aware
rank truncation acts as denoising by removing weight directions that cause
overconfident, incorrect predictions.

We further show the fitted factors tolerate int8 quantization losslessly
($3.0\times$ per-matrix storage at $+3.80\%$), but that low-rank compression
alone cannot compete with whole-model quantization for storage; we characterize
where the combination of the two is the productive direction.
\end{abstract}

\section{Introduction}

A pretrained transformer's dense layers are highly overparameterized by
spectral criteria: singular value spectra of attention projections decay
fast enough that a third of the rank captures most of the Frobenius mass.
Yet replacing these matrices by their truncated SVD collapses full-stack
language modeling performance --- in our measurements, $+21\%$ perplexity on
GPT-2 Small and $+72\%$ on Gemma-3-1B at one-sided $3\times$ compression,
even though per-layer output error under teacher inputs is small.

This failure has been attributed to two competing mechanisms:
\begin{enumerate}[label=(\arabic*)]
  \item \textbf{Drift limit}: approximation errors compound through the
    residual stream; later layers receive off-manifold inputs they were
    never trained on, so layerwise equivalence does not imply network
    equivalence.
  \item \textbf{Norm mismatch}: the weight-space Frobenius norm in which SVD
    is optimal is the wrong norm --- what matters is the error under the
    activation distribution, and SVD wastes rank on directions that never
    fire.
\end{enumerate}
We pre-registered a decision rule to distinguish these at 7B scale, and we
find the answer is overwhelmingly (2): a closed-form fit that minimizes
expected activation-space error recovers a full compression tier that plain
SVD destroys. The ``drift wall'' observed at $3\times$ compression across
model scales ($+21.4\%$ Small, $+9.9\%$ Medium, $+7.5\%$ Large in our
earlier phases) is therefore largely an artifact of the approximation
method, not a fundamental limit.

\paragraph{Contributions.}
\begin{enumerate}[leftmargin=1.5em]
  \item \textbf{A closed-form activation-weighted low-rank fit} with a
    weight-space anchor (Section~\ref{sec:method}):
    $M^{*} = C(G+\beta n I)^{-1}$, truncated to rank $r$ in the
    $G$-weighted norm. No gradient steps; per-matrix solve in seconds.
  \item \textbf{Cross-scale validation} (Section~\ref{sec:results}):
    Testing on GPT-2 Small (124M), Gemma-3-1B, and Qwen2.5-7B against
    matched SVD controls. The weighted fit is $5$--$7\times$ better in
    perplexity delta at every tested scale.
  \item \textbf{Evidence that compression can improve on the teacher}
    (Section~\ref{sec:denoise}): The compressed Gemma-3-1B achieves
    $-4.45\%$ perplexity (and $-7.34\%$ on 200 texts) below the teacher
    with no fine-tuning. This is possible because activation-weighted
    truncation removes noisy weight directions that cause overconfident,
    incorrect predictions. The compressed model shows higher entropy (less
    overconfidence) and lower NLL (better calibration).
  \item \textbf{Honest storage accounting} (Section~\ref{sec:storage}):
    Factor quantization to int8 is lossless ($3.0\times$ per matrix), but
    whole-model quantization dominates pure storage. The productive
    direction is combining low-rank structure with quantized
    representations.
  \item \textbf{Negative results and pitfalls} (Appendix~\ref{app:pitfalls}):
    Drift-aware refitting on student inputs does not help; Conv1D bias
    double-counting and alternating subspace iteration divergence both cause
    order-of-magnitude blowups.
\end{enumerate}

\section{Related Work}

\paragraph{Quantization with activation awareness.} OBQ/OBS and the
GPTQ/SparseGPT line~\cite{frantar2022gptq,frantar2023sparsegpt} established
that weight-space rounding error should be
compensated under the second-order (activation Gram / Hessian) statistic of
calibration inputs. AWQ~\cite{lin2023awq} protects salient channels by
activation magnitude.
Our objective is the low-rank analogue of this principle.

\paragraph{Activation-aware low-rank compression.} ASVD~\cite{yuan2023asvd}
reparameterizes weights by activation statistics before
SVD and handles outlier channels, demonstrating that activation awareness
is the lever for LLM low-rank compression. Our closed-form anchored fit is
an independent derivation in the same family; to our knowledge the
beats-the-teacher denoising effect (Section~\ref{sec:denoise}) has not been
reported for training-free rank truncation. EoRA~\cite{liu2024eora}
compensates low-rank error with an added eigenspace term;
QEP~\cite{arai2025qep} patches cross-layer error propagation for
quantization.

\paragraph{Hybrid low-rank + quantization.} SVDQuant~\cite{li2024svdquant}
migrates outliers into an fp16 low-rank branch and quantizes the residual
to 4 bits. Our factor-quantization result (Section~\ref{sec:storage}) shows
the low-rank branch itself survives int8 losslessly, reinforcing the hybrid
design space.

\paragraph{Classical weighted low-rank approximation.} The problem
$\min_{\operatorname{rank}(\Wh)\le r} \mathbb{E}\|(W-\Wh)X\|$ is a known
non-convex problem; our contribution is the specific anchored closed-form +
weighted-truncation pipeline and its systematic validation on modern LLMs
against matched controls.

\section{Method}\label{sec:method}

\subsection{Problem setting}
For a linear module with weight $W \in \mathbb{R}^{d_{out}\times d_{in}}$
(column-space convention $y = Wx$; for GPT-2 Conv1D modules $W$ is
transposed accordingly), we observe calibration input activations
$x_1,\dots,x_n \in \mathbb{R}^{d_{in}}$ captured by a forward hook, and
seek a rank-$r$ approximation $\Wh$ minimizing
\begin{equation}\label{eq:obj}
  \mathcal{L}(\Wh) = \frac{1}{n}\sum_{i=1}^{n} \|(W-\Wh)x_i\|^2
                     + \beta \|\Wh - W\|_F^2 .
\end{equation}
The first term is the activation-space error (what the network actually
sees); the second is a weight-space anchor of strength $\beta \ge 0$. The
anchor is mandatory, not optional: with $\beta = 0$ the problem is
ill-posed in directions the calibration data never excites, and empirical
fits diverge there.

\subsection{Closed-form solution}
Let $G = XX^\top \in \mathbb{R}^{d_{in}\times d_{in}}$ be the activation
Gram matrix and $C = (Y-b)X^\top + \beta n W$ the cross-moment (with bias
$b$ subtracted from captured outputs where present --- see
Appendix~\ref{app:pitfalls}). The unconstrained minimizer of
\eqref{eq:obj} is
\begin{equation}\label{eq:mstar}
  M^{*} = C\,(G + \beta n I)^{-1},
\end{equation}
computed in the eigenspace of $G$ with a relative eigenvalue floor
$\varepsilon_{rel} = 10^{-3}$ and Tikhonov ridge
$\lambda = 0.01 \cdot \overline{\operatorname{diag}(G_k)}$ for numerical
safety.

\subsection{Rank-$r$ truncation in the correct norm}
Truncating $M^{*}$ by plain SVD would again optimize the wrong norm. We
truncate in the $G$-weighted norm --- the norm in which \eqref{eq:obj}
measures error --- via the SVD of
\begin{equation}\label{eq:trunc}
  Z = M^{*}(G+\beta n I)^{1/2}, \qquad Z = U_z \Sigma V_z^\top,
\end{equation}
and set $\Wh = U_z[:,{:}r]\, U_z[:,{:}r]^\top M^{*}$. This is the exact
rank-$r$ minimizer of \eqref{eq:obj} given the full-rank optimum $M^{*}$;
the per-matrix solve is closed-form and takes seconds on a single GPU.

\subsection{Calibration and hyperparameters}
Calibration: 16 chunks of 512 tokens from the WikiText-2~\cite{merity2016} \emph{train}
split (never the test split), teacher-forced through the untouched model
($\alpha = 0$: teacher inputs; student-input refitting was tested and is
inferior, Section~\ref{sec:ablations}). One hyperparameter $\beta$:
optimal in $[0.1, 0.3]$; all reported results use $\beta = 0.1$. Rank
convention: one-sided, $r = \min(d_{in}, d_{out})/k$.

\section{Experimental Setup}

\paragraph{Models.} GPT-2 Small (124M), \texttt{google/gemma-3-1b-it},
\texttt{Qwen/Qwen2.5-7B-Instruct} (rev \texttt{a09a3545}). Targets:
attention output projections in every layer (\texttt{attn.c\_proj},
\texttt{self\_attn.o\_proj}) --- the family identified as compressible in
Section~\ref{sec:families}.

\paragraph{Protocol (identical across all runs).} Calibration: WikiText-2
train, $16\times512$ tokens (the GPT-2 sweep used $64\times512$).
Evaluation: WikiText-2 test, 50 texts $>$50 chars truncated to 256 tokens
(Gemma additionally verified on 200 texts). Metric: perplexity delta vs the
untouched fp-matched baseline. Control: plain truncated SVD at the
identical rank --- same storage, same modules, same eval. All runs on a
single RTX 4070 Ti 12\,GB; Qwen-7B in fp16 with partial CPU offload
(weights read from safetensors shards, written via
\texttt{set\_module\_tensor\_to\_device}).

\paragraph{Pre-registered decision rule (7B validation).} Weighted-fit
delta $\le 10\%$ supports the redundancy hypothesis (compression scales);
$\ge 20\%$ falsifies it; between is inconclusive.

\section{Results}\label{sec:results}

\subsection{Cross-scale main result (one-sided $3\times$, all output projections)}

\begin{table}[h]
\centering
\begin{tabular}{lrrr}
\toprule
Model & Baseline PPL & Plain SVD & Weighted fit ($\beta=0.1$) \\
\midrule
GPT-2 Small (124M)      & 56.47 & $+21.34\%$ & $\mathbf{+3.58\%}$ \\
Gemma-3-1B              & 70.40 & $+71.86\%$ & $\mathbf{-4.45\%}$ \\
Qwen2.5-7B-Instruct     & 12.29 & $+26.31\%$ & $\mathbf{+3.84\%}$ \\
\bottomrule
\end{tabular}
\caption{Held-out perplexity delta at one-sided $3\times$ rank on all
attention output projections. Weighted-vs-SVD gap: $5.9\times$,
$16.2\times$, $6.9\times$.}
\label{tab:main}
\end{table}

The pre-registered 7B gate returns \textbf{SUPPORTS}: the earlier
``failure wall'' at $3\times$ was a property of the approximation norm, not
of scale.

\subsection{Ratio frontier (GPT-2 Small, c\_proj $\times$12)}

\begin{table}[h]
\centering
\begin{tabular}{lrrrrr}
\toprule
& $2\times$ (r=384) & $3\times$ (r=256) & $4\times$ (r=192) & $6\times$ (r=128) & $8\times$ (r=96) \\
\midrule
SVD      & $+5.50\%$  & $+21.34\%$ & $+152.13\%$ & $+261.86\%$ & $+374.49\%$ \\
Weighted & $\mathbf{+0.10\%}$ & $\mathbf{+3.58\%}$ & $\mathbf{+17.94\%}$ & $+53.38\%$ & $+94.18\%$ \\
\bottomrule
\end{tabular}
\caption{Perplexity delta across compression ratios. The weighted fit
recovers a full compression tier: weighted $4\times$ beats SVD $3\times$ in
PPL while storing less.}
\label{tab:frontier}
\end{table}

\subsection{Ablations}\label{sec:ablations}

\paragraph{$\beta$ sweep.} Optimum $\beta \in [0.1, 0.3]$; $\beta=0$
diverges in unobserved directions.

\paragraph{Drift-awareness (student inputs).} Fitting against the
\emph{student's} own activations ($\alpha=1$, sequential refit) loses
$\sim$1\,pp to teacher-input fitting ($\alpha=0$) at every matched $\beta$.
Anticipating drift is not the lever; fixing the norm is.

\paragraph{Robustness.} Disjoint calibration (chunks 64--127 instead of
0--63): $+4.94\%$ vs $+3.58\%$ --- mild sensitivity to calibration data,
no collapse.

\subsection{Matrix-family study (GPT-2 Small, one-sided $3\times$)}\label{sec:families}

\begin{table}[h]
\centering
\begin{tabular}{lrr}
\toprule
Family & SVD & Weighted \\
\midrule
attn.c\_proj & $+21.34\%$  & $\mathbf{+3.58\%}$ \\
attn.c\_attn & $+807.70\%$ & $+37.17\%$ \\
attn (full)  & $+1072.37\%$& $+49.11\%$ \\
mlp.c\_proj  & $+546.12\%$ & $+76.65\%$ \\
mlp.c\_fc    & $+7544.82\%$& $+2149.19\%$ \\
\bottomrule
\end{tabular}
\caption{Weighted fitting wins in every family by $5$--$25\times$, but MLP
matrices retain no usable redundancy at $3\times$ at these sizes.
Compressible family = attention output projections.}
\label{tab:families}
\end{table}

\subsection{Factor quantization (Qwen2.5-7B, rank 1194)}

\begin{table}[h]
\centering
\begin{tabular}{lrr}
\toprule
Factor precision & Per-matrix storage & PPL delta \\
\midrule
fp16                    & $1.5\times$ & $+3.86\%$ \\
int8 (per-row symmetric)& $3.0\times$ & $\mathbf{+3.80\%}$ \\
int4                    & $6.0\times$ & $+8.22\%$ \\
\bottomrule
\end{tabular}
\caption{int8 factors are lossless --- the fitted subspace is robust to
8-bit factor noise. int4 costs $+4.4$\,pp but remains inside the
$\le 10\%$ gate.}
\label{tab:factorquant}
\end{table}

\subsection{Hybrid residual split (Qwen2.5-7B, SVDQuant-style)}

Per-matrix hybrid: $\hat W = UV + Q_{\mathrm{int4}}(W - UV)$ --- a small
low-rank branch (rank $r=256$, fp16 factors) absorbs structure the grid
cannot represent; the residual is group-wise symmetric int4 (group $128$).
Applied to all 28 \texttt{o\_proj}, same run contract.

\begin{table}[h]
\centering
\begin{tabular}{lccc}
\toprule
Variant & PPL & $\Delta$ & storage/matrix \\
\midrule
Baseline                        & 12.29 & ---            & 25.69\,MB \\
Pure int4 (reference)           & 12.39 & $+0.78\%$      & 6.62\,MB ($3.88\times$) \\
Hybrid, SVD branch              & 12.37 & $+0.60\%$      & 10.29\,MB ($2.50\times$) \\
Hybrid, weighted branch         & \textbf{12.36} & $\mathbf{+0.56\%}$ & 10.29\,MB ($2.50\times$) \\
\bottomrule
\end{tabular}
\caption{Both pre-registered checks pass: the weighted branch beats the
matched-storage SVD branch and beats the cheaper pure-int4 reference. The
margin on \texttt{o\_proj} is small (+0.22\,pp) because int4 alone is
nearly free on this family --- \texttt{o\_proj} is easy for both low-rank
and quantization. The hybrid's payoff should concentrate on outlier-heavy
families (MLPs, 80.7\% of parameters), where plain low-rank fails
(Table~\ref{tab:families}); this motivates per-family routing
(Section~\ref{sec:storage}).}
\label{tab:hybrid}
\end{table}

\section{Analysis: Compression as Denoising}\label{sec:denoise}

\subsection{Why Compression Can Improve on the Teacher}

A compressed model outperforming its teacher is counterintuitive. Normally,
we expect compression to degrade performance. However, our results show that
activation-weighted low-rank fitting can actually \textbf{improve} the
model's language modeling ability.

\paragraph{The key insight.} Transformer weight matrices contain directions
that are never excited by natural text. These ``dead directions'' carry no
signal but contribute to weight-space capacity for spurious logits ---
incorrect predictions. By truncating these directions through
activation-weighted fitting, we remove noise from the model.

\paragraph{Evidence from Gemma-3-1B:}
\begin{itemize}
  \item Teacher perplexity: 70.40
  \item Compressed model perplexity: 67.47 ($-4.45\%$)
  \item On 200 texts: $-7.34\%$ improvement
\end{itemize}

The compressed model is \textbf{not} a copy of the teacher. Logit-space
comparison shows:
\begin{itemize}
  \item Top-1 agreement: only 74.4\% (the models disagree on 1 in 4 tokens)
  \item Logit cosine: 0.9925 (high similarity but not identical)
  \item KL divergence: 0.268 / 0.292 (substantial difference)
\end{itemize}

\paragraph{The denoising mechanism:}
\begin{enumerate}
  \item Natural text only excites certain directions in weight matrices
  \item Unused directions contribute to overconfident, incorrect predictions
  \item Activation-weighted truncation removes these noise directions
  \item The model becomes less overconfident (entropy increases 12\%)
  \item Better calibration leads to lower perplexity on held-out text
\end{enumerate}

This is analogous to regularization in machine learning --- removing noisy
parameters can improve generalization. The difference is that our method
identifies ``noise'' directions using activation statistics, not weight
magnitude alone.

\subsection{Teacher--Student Divergence}

The Gemma result ($-4.45\%$; $-7.34\%$ on 200 texts) means the compressed
model is a \emph{better} language model on held-out text than the teacher,
with no fine-tuning. Logit-space comparison on 30 texts ($\sim$5.9k tokens)
shows the student is not a near-copy:

\begin{table}[h]
\centering
\begin{tabular}{lr}
\toprule
Metric & Value \\
\midrule
KL(orig$\|$sub) / KL(sub$\|$orig) & 0.268 / 0.292 \\
Logit cosine                      & 0.9925 \\
Top-1 / top-5 agreement           & 74.4\% / 96.3\% \\
Softmax entropy (orig $\to$ sub)  & 2.19 $\to$ 2.47 \\
\bottomrule
\end{tabular}
\caption{Teacher--student divergence of the Gemma compressed model.}
\label{tab:divergence}
\end{table}

Entropy rises 12\% while NLL falls: the student is less overconfident and
better calibrated. Interpretation: directions of $W$ that are never excited
by natural text carry no signal but contribute weight-space capacity for
spurious logits; activation-weighted truncation removes them, flattening
spurious peaks and moving mass onto correct tokens --- denoising.

At 7B, qualitative generation on three solvable prompts (arithmetic,
factual QA, code) shows the int8-factor model tracks the baseline nearly
word-for-word, the int4 model paraphrases with correct content, and even
the $+26\%$-PPL SVD control remains coherent on easy tasks --- greedy
generation on short prompts is too insensitive a probe for this kind of
damage, confirming held-out PPL / logit divergence as the right primary
metrics.

\gap\ The denoising claim currently rests on one architecture, one family,
and 1--2 eval sizes. Required before claiming it: 3 seeds, replication on a
second architecture, pre-registered calibration/eval protocol.

\section{Storage Accounting vs Quantization}\label{sec:storage}

Factored storage of a rank-$d/3$ approximation of a $d\times d$ matrix is
$d^2 \cdot 2/3$ parameters $\to$ $1.5\times$; ``$k\times$ rank'' is not
``$k\times$ storage'' ($k/2\times$ for square matrices). On Qwen2.5-7B
(o\_proj = 5.09\% of params):

\begin{table}[h]
\centering
\begin{tabular}{lrrl}
\toprule
Method & Storage & PPL cost & Scope \\
\midrule
Weighted $3\times$-rank, int8 factors & $3.0\times$/matrix $\to$ $\sim$3.4\% whole model & $+3.80\%$ & o\_proj \\
int8 quantization                     & $2.0\times$ whole model & $\sim$0--0.1\% & all weights \\
int4 (GPTQ/AWQ)                       & $\sim$3.6$\times$ whole model & $\sim$0.3--1\% & all weights \\
\bottomrule
\end{tabular}
\caption{Storage accounting. Low-rank compression alone cannot compete with
quantization for storage.}
\label{tab:storage}
\end{table}

\paragraph{Verdict.} Low-rank compression alone cannot compete with
quantization for storage. Its value is quality-at-rank and the denoising
effect; the productive frontier is combination --- now measured
(Table~\ref{tab:hybrid}): the weighted-fit branch beats both the SVD branch
and pure int4 on \texttt{o\_proj}. The design space that remains open is
per-family routing --- hybrid residual split on outlier-heavy MLPs (80.7\%
of parameters), low-rank where it wins, plain quantization where int4 is
free.

\subsection{What is actually stored: a byte-level worked example}\label{sec:bytes}

We make the conversion concrete for one Qwen2.5-7B o\_proj matrix
($d_{out} = d_{in} = 3584$, rank $r = 1194$, the measured configuration of
Table~\ref{tab:factorquant}).

\paragraph{Before conversion (on disk today).} The matrix is stored dense:
$3584 \times 3584 = 12{,}845{,}056$ weights, 2 bytes each in fp16
$= 25{,}690{,}112$ bytes $= 25.69$\,MB. Every weight is stored regardless of
whether any real input ever excites its direction.

\paragraph{The conversion.} Calibration activations $X$ are collected; we
solve \eqref{eq:mstar}--\eqref{eq:trunc} and obtain factors
$U \in \mathbb{R}^{3584\times1194}$ and $V \in \mathbb{R}^{1194\times3584}$
with $\Wh = UV$. The dense matrix is then \emph{discarded}; only the
factors are serialized. The inference-time computation $y = Wx$ becomes
$y = U(Vx)$: one $1194\times3584$ matmul followed by one
$3584\times1194$ matmul.

\paragraph{After conversion, three precisions (per matrix).}
\begin{table}[h]
\centering
\begin{tabular}{lrrr}
\toprule
Serialization & Codes / weights & Overhead & Total \\
\midrule
fp16 factors  & $8{,}558{,}592\ \text{params} \times 2$\,B (two $3584\times1194$ factors) & --- & 17.12\,MB ($1.50\times$) \\
int8 factors  & $8{,}558{,}592\ \text{codes} \times 1$\,B & $(3584+1194)$ scales $\times 2$\,B $= 9.6$\,KB & 8.57\,MB ($3.00\times$) \\
int4 factors  & $8{,}558{,}592\ \text{codes} \times 0.5$\,B & 9.6\,KB & 4.29\,MB ($5.99\times$) \\
\bottomrule
\end{tabular}
\caption{Exact bytes stored for one Qwen2.5-7B o\_proj after conversion.
Measured PPL deltas at each point: $+3.86\%$ (fp16), $+3.80\%$ (int8),
$+8.22\%$ (int4).}
\label{tab:bytes}
\end{table}

At int8 the matrix shrinks from 25.69\,MB to 8.57\,MB --- one third of the
space --- at \emph{measured} zero additional quality cost ($+3.80\%$ vs
$+3.86\%$). The per-row scale overhead (9.6\,KB, $0.11\%$ of the codes) is
negligible.

\paragraph{Model-level arithmetic.} Qwen2.5-7B has 28 o\_proj matrices
($5.09\%$ of its 7.07B params). Converting all of them to int8 factors:
$28 \times (25.69 - 8.57) = 479$\,MB saved, i.e.\ the fp16 model drops from
14.14\,GB to $\sim$13.66\,GB ($\sim$3.4\% whole-model saving). This is the
honest ceiling of the method at its current scope, and the reason
Section~\ref{sec:storage} positions combination with quantization, not
competition with it, as the productive frontier.

\section{Limitations}

\begin{enumerate}[leftmargin=1.5em]
  \item \textbf{Scope of compression.} Only attention output projections
    are viable at $3\times$; MLP families break even weighted. Whole-model
    compression requires budget reallocation across families, which we have
    not demonstrated.
  \item \textbf{Seeds.} Single calibration/fit per result \gap; the
    $\pm 1$\,pp robustness run suggests stability but multi-seed statistics
    are missing.
  \item \textbf{Eval sensitivity.} Perplexity on $\le 200$ WikiText-2
    texts; no downstream-task evaluation, no pre-registered drift metric
    alongside PPL \gap.
  \item \textbf{Latency/size accounting.} Factored form changes the compute
    graph; warmed-up latency and peak memory are unmeasured \gap. Without a
    fused low-rank kernel the storage win does not automatically become a
    speed win.
  \item \textbf{Instruct-tuned 7B baseline.} Qwen experiments use the
    Instruct variant (matched to prior cross-architecture protocol);
    base-model PPL might differ in absolute terms.
\end{enumerate}

\section{Revised conclusion}

The current evidence supports activation-weighted low-rank fitting as a
promising empirical method for Gemma attention output projections, with a
clear advantage over plain SVD at the tested rank. It does not yet establish
algorithmic novelty, cross-architecture generality, or deployment success.
Those claims require matched ASVD/EoRA controls, full drift and downstream
evaluation, and a genuinely packed artifact. Direct high-sparsity pruning is
reported as a negative result.

\section{Previous conclusion (superseded)}

The ``drift wall'' of full-stack low-rank compression is largely an
artifact of approximating in the wrong norm. A closed-form
activation-weighted fit with a weight-space anchor, truncated in the
activation Gram norm, is training-free, cheap, and defeats truncated SVD by
$5$--$7\times$ across three architectures from 124M to 7B parameters,
recovers a full compression tier, and --- on Gemma-3-1B --- produces a
compressed model that outperforms its teacher on held-out perplexity, with
a measurable calibration signature. Pure storage compression remains
quantization's territory; the compelling direction is combining data-aware
low-rank structure with quantized representations.

\appendix

\section{Pitfalls with Quantitative Evidence}\label{app:pitfalls}

\paragraph{A.1 Bias double-counting.} Modules whose outputs include a bias
(GPT-2 Conv1D) produce captured $Y = WX + b$; fitting $Y$ while the module
keeps $b$ double-counts it. Layer-0 sanity check showed
$\|WX - Y\|/\|Y\| = 0.216$ --- a 22\% phantom target. Full-stack
consequence: $+3470\%$ PPL until fixed.

\paragraph{A.2 Alternating subspace iteration diverges.} Refitting
$Z = MPG^{1/2}$ against the current solution's energy (ignoring the
cross-moment $C$) diverges: $\|M-W\|_2$ grew $9.9 \to 23.6$ over 4
iterations. Closed-form \eqref{eq:mstar}+\eqref{eq:trunc} is required.

\paragraph{A.3 Offloaded weights are meta.} With accelerate
\texttt{device\_map="auto"}, CPU-offloaded parameters report
\texttt{device == meta}; direct \texttt{.data} access raises. Read from
safetensors shards via the index \texttt{weight\_map}, write with
\texttt{set\_module\_tensor\_to\_device}.

\paragraph{A.4 Student-input refitting is inferior.} Anticipating drift
($\alpha = 1$) loses $\sim$1\,pp to teacher-input fitting at every matched
$\beta$; the sequential refit adds complexity without payoff.

\section{Reproducibility}

Every number above is produced by a standalone script in \texttt{src/} with
output in \texttt{results/}: \texttt{drift\_aware\_svd.py} (+ sweeps),
\texttt{weighted\_one\_family.py}, \texttt{weighted\_gemma\_oproj.py}
(+ \texttt{\_200}), \texttt{weighted\_ratio\_frontier.py},
\texttt{drift\_aware\_robustness.py}, \texttt{weighted\_qwen7b\_oproj.py},
\texttt{weighted\_qwen7b\_factor\_quant.py},
\texttt{hybrid\_residual\_split.py}, \texttt{compare\_gemma\_output.py},
\texttt{compare\_qwen7b\_output.py},
\texttt{storage\_vs\_quantization.py}. The run contract (model revisions,
splits, seeds, budgets) is recorded in the header of each script and
summarized in \texttt{THREAD.md}. Hardware: RTX 4070 Ti 12\,GB, Python
3.14, torch 2.13.0+cu132.

\section*{Pre-submission checklist (\gap\ items)}
\begin{itemize}[leftmargin=1.5em]
  \item 3 seeds for the denoising claim (Gemma $-4.45\%$) and a second
    architecture reproducing improvement
  \item Pre-registered drift metric (hidden-state cosine per block)
    alongside PPL
  \item MLP budget reallocation experiment (adaptive rank across families)
  \item Warmed-up latency + peak-memory accounting for factored inference
  \item Hybrid residual split on MLP families (where the payoff should
    live) and the full per-family routed recipe
  \item Head-to-head against ASVD-style activation reparameterization at
    matched budget (the sharpest novelty test)
  \item Base (non-Instruct) 7B confirmation run
\end{itemize}

\begin{thebibliography}{9}

\bibitem{frantar2022gptq}
E.~Frantar, S.~Ashkboos, T.~Hoefler, and D.~Alistarh.
GPTQ: Accurate post-training quantization for generative pre-trained
transformers.
\emph{arXiv:2210.17323}, 2022.

\bibitem{frantar2023sparsegpt}
E.~Frantar and D.~Alistarh.
SparseGPT: Massive language models can be accurately pruned in one-shot.
\emph{arXiv:2301.00774}, 2023.

\bibitem{lin2023awq}
J.~Lin, J.~Tang, H.~Tang, S.~Yang, W.-M.~Chen, W.-C.~Wang, G.~Xiao,
X.~Dang, C.~Gan, and S.~Han.
AWQ: Activation-aware weight quantization for LLM compression and
acceleration.
\emph{arXiv:2306.00978}, 2023.

\bibitem{yuan2023asvd}
Z.~Yuan, Y.~Shang, Y.~Song, D.~Yang, Q.~Wu, Y.~Yan, and G.~Sun.
ASVD: Activation-aware singular value decomposition for large language
model compression.
\emph{arXiv:2312.05821}, 2023.

\bibitem{liu2024eora}
S.-Y.~Liu, M.~Khadkevich, N.~C.~Fung, C.~Sakr, C.-H.~H.~Yang, C.-Y.~Wang,
S.~Muralidharan, H.~Yin, K.-T.~Cheng, J.~Kautz, Y.-C.~F.~Wang,
P.~Molchanov, and M.-H.~Chen.
EoRA: Fine-tuning-free compensation for compressed LLM with eigenspace
low-rank approximation.
\emph{arXiv:2410.21271}, 2024.

\bibitem{arai2025qep}
Y.~Arai and Y.~Ichikawa.
Quantization error propagation: Revisiting layer-wise post-training
quantization.
\emph{arXiv:2504.09629}, 2025.

\bibitem{li2024svdquant}
M.~Li, Y.~Lin, Z.~Zhang, T.~Cai, X.~Li, J.~Guo, E.~Xie, C.~Meng,
J.-Y.~Zhu, and S.~Han.
SVDQuant: Absorbing outliers by low-rank components for 4-bit diffusion
models.
\emph{arXiv:2411.05007}, 2024.

\bibitem{merity2016}
S.~Merity, C.~Xiong, J.~Bradbury, and R.~Socher.
Pointer sentinel mixture models.
\emph{arXiv:1609.07843}, 2016.

\end{thebibliography}

\end{document}
